\chapter{把工艺数据变成下一版设计：一个可运行的定量闭环}

\begin{chaptergoal}
把前面章节里的 film map、CD map、resonance map 和 DfM 连接成一个可以真正执行的数据流程：先定义可追溯的数据表，再建立一个透明的小信号 surrogate，检查它能解释多少 frequency scatter，最后只对可重复的 systematic component 做下一版设计补偿；同时理解为什么“相关”不等于“工艺原因已经证明”。
\end{chaptergoal}

到 v0.1 为止，我们已经知道应该测什么、记录什么、怎样避免把一个异常 notch 直接归因到某一道工艺。v0.2 往前再走一步：\textbf{把这些记录变成能跑的模型。}

这里故意先使用一个很小的 synthetic example，而不是一上来做复杂机器学习。原因很简单：如果一个 64-resonator toy model 都不能把 ID、单位、sensitivity、collision rule 和 correction 写清楚，那么换成几千个像素只会把混乱放大。

\processchain{wafer metrology map $\rightarrow$ per-resonator predictor $\rightarrow$ measured frequency residual $\rightarrow$ collision audit $\rightarrow$ repeatability check $\rightarrow$ next-GDS correction}

\section{第一步：所有数据必须先通过 resonator ID 对齐}

对第 $i$ 个 resonator，至少需要把下面几类量放到同一个可连接的数据系统中：

\begin{itemize}
  \item 设计侧：$f_{i}^{\mathrm{target}}$、GDS revision、物理坐标；
  \item 薄膜侧：$R_{\square,i}$、厚度或其他可空间插值的 material proxy；
  \item 几何侧：linewidth、gap、coupler 等 measured CD；
  \item 低温侧：$f_{0,i}^{\mathrm{meas}}$、$Q_{r,i}$、$Q_{i,i}$、$Q_{c,i}$；
  \item 追溯侧：batch / wafer / die / package / cooldown ID。
\end{itemize}

因此 v0.2 的第一件事不是“拟合一个模型”，而是先让
\[
\text{wafer coordinate}
\leftrightarrow
\text{physical pixel}
\leftrightarrow
\text{resonator ID}
\leftrightarrow
\text{measured notch}
\]
能够稳定连接。

本仓库从 v0.2 开始用 \texttt{notes/DATA\_CONTRACT.md} 固定这一层。对应的两个核心表是：

\begin{itemize}
  \item \texttt{wafer\_metrology\_map.csv}：保存物理位置对应的 film / CD / defect 数据；
  \item \texttt{resonator\_map.csv}：连接 design frequency、physical pixel 与低温 fit 结果。
\end{itemize}

\section{第二步：先定义我们到底在拟合什么}

最直接的目标量是 fractional frequency error：
\[
\epsilon_i
=
\frac{
f_{0,i}^{\mathrm{meas}}-f_{i}^{\mathrm{target}}
}{
f_{i}^{\mathrm{target}}
}.
\]

如果只看 $\epsilon_i$ 的 histogram，我们只能知道 spread 多大；如果把它与同位置的 film / geometry 数据对齐，就可以尝试建立 predictor。

对薄膜部分，第 3 章已经给出近似关系
\[
\frac{\delta f_0}{f_0}
\approx
-\frac{\alpha}{2}
\frac{\delta L_k}{L_k}.
\]

在一个只用于建立工程直觉的近似里，如果 $T_c$ 变化暂时忽略，并把 sheet kinetic inductance 的变化主要接到 $R_\square$，可以写成
\[
\epsilon_{R_\square,i}
\approx
-\frac{\alpha}{2}
\frac{R_{\square,i}-R_{\square,0}}{R_{\square,0}}.
\]

这部分至少有明确的物理方向。

几何部分则更适合写成 sensitivity：
\[
s_w
=
\frac{\partial \ln f_0}{\partial \ln w},
\qquad
\epsilon_{w,i}
\approx
s_w
\frac{w_i-w_0}{w_0}.
\]

这里的 $s_w$ \textbf{不是材料常数，也不是本讲义可以替实验室替你规定的数字}。它应该来自 Sonnet/CST/解析 LC sensitivity，或者来自已经验证过的 measured-data regression。

于是一个最小 predictor 可以写成
\[
\hat{\epsilon}_i
=
-\frac{\alpha}{2}
\frac{R_{\square,i}-R_{\square,0}}{R_{\square,0}}
+
s_w
\frac{w_i-w_0}{w_0}.
\]

真正值得看的不是 $\hat{\epsilon}_i$ 本身，而是 residual：
\[
r_i
=
\epsilon_i-\hat{\epsilon}_i.
\]

如果 residual 仍然和 wafer position、gap、package revision 或某个 process split 强相关，就说明模型还漏掉了结构。

\section{第三步：用一个 synthetic 64-resonator 数据集先把逻辑跑通}

仓库中的
\[
\texttt{examples/wafer\_to\_resonance\_demo.py}
\]
使用一个固定随机种子的 $8\times8$ synthetic array。它故意包含：

\begin{itemize}
  \item 一个带有空间梯度的 $R_\square$ map；
  \item 一个带有空间梯度的 linewidth map；
  \item 一个额外的小随机 residual；
  \item 64 个 nominally 等间隔的 design frequencies；
  \item 一个显式的 loaded-$Q$ 和 collision rule。
\end{itemize}

这个例子中的数值全部是教学数据，不代表任何真实 Al LEKID 工艺能力，也不应该被复制成 fabrication specification。

运行：

\begin{verbatim}
cd volume2-fabrication
python3 examples/wafer_to_resonance_demo.py
\end{verbatim}

脚本会生成：

\begin{itemize}
  \item \texttt{wafer\_metrology\_map.csv}；
  \item \texttt{resonator\_map.csv}；
  \item \texttt{summary.txt}。
\end{itemize}

固定 seed 下，这个 toy model 在补偿前得到大约
\[
\sigma_{\epsilon}\approx 860~\mathrm{ppm},
\]
模型补偿后剩余大约
\[
\sigma_{r}\approx 111~\mathrm{ppm}.
\]

同一个例子中，在“相邻实测 resonance 间隔小于 5 个 loaded linewidth”这个\textbf{仅用于演示}的判据下，collision pair 从 2 组降到 0 组。

\begin{warningbox}
这里最重要的不是 ``860 ppm''、``111 ppm'' 或 ``5 linewidths'' 这些数字。它们只是 deterministic regression target，方便 CI 检查示例没有被改坏。真正实验中，sensitivity、process distribution 和 collision margin 都必须重新定义。
\end{warningbox}

\section{第四步：为什么可以做 pre-compensation}

如果上一批已经得到一个可重复的 predictor $\hat{\epsilon}_i$，而下一批仍使用足够相似的 process，可以把目标频率 $f_i^{\mathrm{target}}$ 反推成新的 drawn/design frequency：

\[
f_i^{\mathrm{draw,new}}
=
\frac{
f_i^{\mathrm{target}}
}{
1+\hat{\epsilon}_i
}.
\]

直觉上，这是在版图里提前抵消一个已经被测量证明存在、而且能够重复出现的 systematic shift。

但这一步有一个非常严格的前提：

\begin{center}
\textbf{只能补偿“可重复的偏差”，不能用固定 correction 去消灭不可重复的随机 scatter。}
\end{center}

如果某一批的 wafer gradient 下一批就换了方向，那么把它写回 GDS 反而会制造新的系统误差。

\section{第五步：collision audit 要把 linewidth 也放进去}

只比较 nominal frequency spacing 还不够。一个简单的工程判据可以写成

\[
|f_i-f_j|
<
m\,
\max
\left(
\frac{f_i}{Q_{r,i}},
\frac{f_j}{Q_{r,j}}
\right),
\]

满足时先标记为 collision candidate。

这里 $m$ 是 readout / fitting / tone-placement 所需要的 margin，不是物理常数。不同 readout architecture、拟合算法和动态范围会允许不同的安全间隔。

因此一个 array-yield model 至少应该同时知道：

\begin{enumerate}
  \item design spacing；
  \item fabrication frequency scatter；
  \item loaded linewidth；
  \item readout margin；
  \item out-of-band / unusable-band 条件。
\end{enumerate}

这比单纯写“两个 resonance 距离小于 1 MHz 就算 collision”更容易迁移到不同系统。

\section{第六步：不要在同一批数据上既拟合又宣布成功}

toy model 可以这样做，因为它只是演示。但真实 fabrication 数据必须避免一个很常见的陷阱：用同一批 wafer 拟合 sensitivity，再用同一批 wafer 证明模型很好。

至少可以采用下面一种做法：

\begin{itemize}
  \item Batch A 拟合 systematic film/CD predictor；
  \item Batch B 不重新调参数，直接检查预测 residual；
  \item 只有当方向、幅度和 spatial structure 都具有 repeatability，才考虑写入下一版 GDS compensation。
\end{itemize}

如果 batch 数还很少，就不要急着叫它“process model”；更准确的名字是 working hypothesis。

\section{第七步：correlation 不是 root cause}

假设你发现
\[
R_\square(x,y)
\]
和
\[
\epsilon(x,y)
\]
看起来高度相关，也不能立刻得出“频率偏差就是方阻造成的”。

可能存在的 confounder 包括：

\begin{itemize}
  \item 膜厚和方阻本来就共变；
  \item deposition uniformity 与 lithography focus 也共享 wafer position；
  \item resonator mapping 有系统错配；
  \item temperature / package / readout calibration 具有位置相关偏差；
  \item 另一个没有记录的 process variable 同时驱动两者。
\end{itemize}

更强的证据来自 controlled split、独立 batch、as-fabricated EM model 和多种 metrology 同时闭合。

\begin{checkpointbox}
模型的目标不是给每个异常找一个“最像的原因”，而是把 hypothesis space 越缩越小，并且让下一批实验能真正区分剩余的几个解释。
\end{checkpointbox}

\section{v0.2 建议形成的四张标准图}

当真实数据开始积累后，每一批至少应该自动生成四类图：

\begin{enumerate}
  \item wafer map：$R_\square$ / thickness / CD；
  \item resonance map：$\epsilon_i$、$Q_i$、$Q_c$；
  \item predictor-vs-measurement：$\hat{\epsilon}_i$ 对 $\epsilon_i$；
  \item residual map：$r_i$ 在 wafer coordinate 上的空间结构。
\end{enumerate}

前三张告诉我们模型解释了什么，第四张告诉我们还漏了什么。

后续 v0.2 会在这个标准数据接口上再加入 plotting 和 measured-data adapter，而不是另起一套互不兼容的 notebook。

\begin{measurebox}
\textbf{一个最小可计算 fabrication-feedback loop：}
统一 ID 与单位；wafer metrology map；resonator map；明确 sensitivity source；predictor；residual；collision rule；跨 batch repeatability；只对稳定 systematic component 做 next-GDS correction。
\end{measurebox}

\section{这一章真正的出口}

读完以后，面对一批新的 KID wafer，不应该再只问“频率偏了多少”。更有用的问题是：

\begin{enumerate}
  \item 哪部分偏差能被 measured film/CD data 解释？
  \item 哪部分仍留在 residual 里？
  \item 这个 predictor 在下一批还能成立吗？
  \item collision 是 nominal spacing 不够，还是 fabrication scatter 太大？
  \item 下一版应该补偿 systematic bias，还是重新设计一个更低 sensitivity 的 geometry？
\end{enumerate}

当这些问题能够用同一套表格和脚本回答时，fabrication learning 才真正进入闭环。
